\documentclass[aspectratio=169]{beamer}
\usetheme{Warsaw}

\title[Shortest Paths]{Lecture 9: Shortest Paths in Weighted Graphs}
\subtitle{CS 341: Algorithms and Data Structures}
\author[E. Marsh]{Dr.\ Elena Marsh}
\institute[Rivergate]{Department of Computer Science, Rivergate University}
\date{Spring Term 2026}

\begin{document}

\begin{frame}
  \titlepage
\end{frame}

\begin{frame}{Today's plan}
  \tableofcontents
\end{frame}

\section{Dijkstra's algorithm}

\begin{frame}{Dijkstra's algorithm}
  Greedy relaxation from a source $s$ over non-negative edge weights.
  \begin{itemize}
    \item Maintain tentative distances $d[v]$, initially $d[s] = 0$ and $d[v] = \infty$.
    \item Repeatedly extract the unvisited vertex with the smallest $d[v]$.
    \item Relax each outgoing edge: $d[v] \leftarrow \min(d[v],\, d[u] + w(u, v))$.
  \end{itemize}
  With a binary heap the running time is $O((n + m)\log n)$.
\end{frame}

\section{Handling negative weights}

\begin{frame}{Bellman-Ford and negative weights}
  When edges may have negative weight, greedy choices fail.
  \begin{block}{Bellman-Ford}
    Relax every edge $n - 1$ times: $O(nm)$ total. A further improving pass
    certifies a negative cycle reachable from $s$.
  \end{block}
  \begin{alertblock}{Common exam mistake}
    Adding a constant to every edge weight does not preserve shortest paths,
    because paths with more edges are penalized more.
  \end{alertblock}
\end{frame}

\begin{frame}{Wrapping up}
  \begin{itemize}
    \item Reading: CLRS chapter 22, sections 1 to 3.
    \item Problem set 5 is due Friday at noon.
    \item Next lecture: all-pairs shortest paths and Floyd-Warshall.
  \end{itemize}
\end{frame}

\end{document}
